\section{Method}
In this section, we describe our approach to adapting retrieval-augmented LLMs, by pre-determining the query complexity and then selecting the most fitting strategies for retrieval-augmented LLMs.

\subsection{Preliminaries}
\label{method:pre}
We begin with preliminaries, formally introducing different strategies of retrieval-augmented LLMs.


\paragraph{Non Retrieval for QA}
Let us first define an LLM as a model $\texttt{LLM}$, which takes a sequence of tokens $\vx = [x_1, x_2, ..., x_n]$ as an input and then generates a sequence of tokens $\vy = [y_1, y_2, ..., y_n]$ as an output, which is formalized as follows: $\vy = \texttt{LLM}(\vx)$. Then, in our problem setup for QA, $\vx$ and $\vy$ become the input query ($\vq$) from the user and the generated answer ($\bar{\va}$) from the LLM, respectively: $\vq = \vx$ and $\bar{\va} = \vy$. Also, subsequently, the most naïve LLM-powered QA model can be represented as follows: $\bar{\va} = \texttt{LLM}(\vq)$. Ideally, $\bar{\va}$ should match the actual correct answer $\va$. This non-retrieval-based QA method is highly efficient and could be a somewhat promising approach to handling easy queries, as the size of LLMs becomes extremely large with its effect on storing a large amount of knowledge. However, this approach is largely problematic on queries that require precise or concurrent knowledge of specific people, events, or any subjects beyond the LLMs' internal knowledge.



\paragraph{Single-step Approach for QA}
To address the aforementioned scenarios where $\texttt{LLM}$ may struggle with queries that are not answerable by $\texttt{LLM}$ itself, we can utilize the external knowledge $\vd$, which includes useful information for queries, retrieved from the external knowledge source $\mathcal{D}$ that could be an encyclopedia (e.g., Wikipedia) consisting of millions of documents. Specifically, to obtain such $\vd$ from $\mathcal{D}$, a specific retrieval model is necessary, which returns documents based on their relevance with the given query. This process can be formulated as follows: $\vd = \texttt{Retriever}(\vq; D)$, where $\texttt{Retriever}$ is the retrieval model, with $\vd \in \mathcal{D}$. Here, we can use any off-the-shelf retriever~\cite{bm25, DPR}. 

After the retrieval step is done, we now have a pair of query $\vq$ and its relevant documents $\vd$. Then, in order to augment LLMs with this retrieved external knowledge, we can incorporate it into the input of LLMs, represented as follows: $\bar{\va} = \texttt{LLM}(\vq, \vd)$. This process allows LLMs to gain access to external information contained in $\vd$, which can provide the supplementary context that the internal knowledge of $\texttt{LLM}$ lacks, which can subsequently improve the accuracy and concurrency of LLMs for QA. 

\paragraph{Multi-step Approach for QA}
Even though the aforementioned single-step approach offers significant improvements over non-retrieval for $\vq$ that requires external knowledge, it encounters notable limitations, particularly when dealing with complex queries that necessitate synthesizing information from multiple source documents and reasoning over them. This is where a multi-step approach and reasoning for QA become essential. 

In this multi-step approach, $\texttt{LLM}$ interacts with \texttt{Retriever} in several rounds, progressively refining its understanding of $\vq$, until it formulates the final answer from findings accumulated across these multiple steps. Specifically, the process begins with the initial query $\vq$, and at every retrieval step $i$, new documents $\vd_i$ are retrieved from $\mathcal{D}$ and then incorporated into the input of LLMs, as follows: $\bar{\va}_i = \texttt{LLM}(\vq, \vd_i, \vc_i)$, where the additional context $\vc_i$ can be composed of previous documents and outcomes $(\vd_1, \vd_2, ..., \vd_{i-1}, \bar{\va}_1, \bar{\va}_2, ..., \bar{\va}_{i-1})$, and $\vd_i = \texttt{Retriever}(\vq, \vc_i; D)$\footnote{It is worth noting that implementations of the LLM and retriever vary across different multi-step retrieval-augmented LLM approaches~\cite{ircot, self-ask, react}; therefore, the context $\vc_i$ may incorporate none, some, or all of the previous documents and answers.}. We would like to note that this iterative, multi-step process enables $\texttt{LLM}$ to construct a more comprehensive and extensive foundation to solve queries effectively, specifically adept at complex multi-hop queries where answers depend on interconnected pieces of information. However, it is important to recognize that this multi-step approach can be resource-intensive due to the repeated accesses to $\texttt{Retriever}$ and $\texttt{LLM}$, which entail substantial computational costs.


\subsection{Adaptive-RAG: Adaptive Retrieval-Augmented Generation}
\label{method:ours}

We now introduce our adaptive retrieval-augmented LLMs, which are built upon three different strategies described in the previous section, and which are designed to select the most suitable strategy according to the complexity of queries.



\paragraph{Adapting Retrieval-Augmented LLMs}
Note that in real-world scenarios, not all $\vq$ from users have the same level of complexity, necessitating tailored strategies for handling each query. In other words, employing the most basic, non-retrieval-based approach $\texttt{LLM}(\vq)$ to respond to the complex query $\vq$ would be also ineffective (Figure~\ref{fig:concept}, A); conversely, using a more elaborate multi-step approach $ \texttt{LLM}(\vq, \vd, \vc)$ for simple $\vq$ would be inefficient (Figure~\ref{fig:concept}, B). Therefore, our adaptive framework is designed to dynamically adjust the query-handling strategy of retrieval-augmented LLMs, which is achieved by determining the complexity of each query before attempting a solution. Notably, this framework can offer a robust middle ground with a range of solutions, from the simplest approach for the most straightforward queries, to the one-step approach for moderate queries, and up to the most comprehensive and rigorous approach for complex queries. In addition, since the operations of $\texttt{LLM}$ and $\texttt{Retriever}$ remain consistent regardless of inputs to them, our method can seeminglessly go back and forth across queries of different complexities, without changing the internal model architecture or parameters during adaption.


\paragraph{Query Complexity Assessment}
To operationalize our adaptive retrieval-augmented LLM framework, we should determine the query complexity, and to achieve this, we propose to model a complexity classifier, whose goal is to return the appropriate complexity level of the given query. Specifically, given the query $\vq$, our classifier can be formulated as follows: $o = \texttt{Classifier}(\vq)$, where $\texttt{Classifier}$ is a smaller Language Model that is trained to classify one of three different complexity levels and $o$ is its corresponding class label. In our classifier design, there are three class labels: `A', `B', and `C', where `A' indicates that $\vq$ is straightforward and answerable by $\texttt{LLM}(\vq)$ itself, `B' indicates that $\vq$ has the moderate complexity where at least a single-step approach $\texttt{LLM}(\vq, \vd)$ is needed, and `C' indicates that $\vq$ is complex, requiring the most extensive solution $\texttt{LLM}(\vq, \vd, \vc)$\footnote{We consider three levels of query complexity, and leave the exploration of more fine-grained complexities as future work.}.



\paragraph{Training Strategy}
The remaining step is to train the smaller Language Model for $\texttt{Classifier}$, to accurately predict its complexity $o$ in response to the given query $\vq$. Yet, there is no annotated dataset available for query-complexity pairs. Hence, we propose to automatically construct the training dataset with two particular strategies. 

To be specific, we first aim at labeling the query complexity based on the results from three different retrieval-augmented LLM strategies, in order to determine the label by its needs. For example, if the simplest non-retrieval-based approach correctly generates the answer, the label for its corresponding query is assigned `A'. Also, to break the tie between different models in providing the label to the query, we provide a higher priority to a simpler model. In other words, if both single-step and multi-step approaches produce the same correct answer while the non-retrieval-based approach fails, we assign label `B' to its corresponding query. 

However, this labeling strategy has a limitation in that not all the queries are assigned labels, since the three retrieval-augmented approaches may all fail to generate the correct answer. On the other hand, the benchmark datasets may already have meaningful inductive biases about the most appropriate retrieval-augmented LLM strategies for their queries, considering the ways they are created (e.g., QA datasets that require sequential reasoning usually necessitate a multi-step approach; while queries of those with labeled single documents can be ideally answerable with the single-step approach). Therefore, for those queries that remain unlabeled after the first labeling step, we assign `B' to queries in single-hop datasets and `C' to queries in multi-hop datasets. Finally, we train $\texttt{Classifier}$ with these automatically-collected query-complexity pairs\footnote{As we automatically assign classifier labels, there might be errors in labeling and might be more advanced strategies to automatically assign labels, which we leave as future work.}, by using a cross-entropy loss. Then, at inference, we can determine the complexity of the query, which is one of \{`A', `B', `C'\}, by forwarding it to $\texttt{Classifier}$: $o = \texttt{Classifier}(\vq)$.
